% =====================================================================
\chapter{Autoevaluación}
\label{ap:autoevaluacion}
% =====================================================================

\section{Puntuación por parte y por módulo}

El examen se corrige con la clave del \cref{tab:clave} y con los criterios entre corchetes de cada solución. El \cref{tab:puntos-parte} recoge la puntuación por parte y el \cref{tab:puntos-modulo}, la puntuación por módulo, que sirve para localizar las debilidades.

\begin{table}[h]
\centering\small
\caption{Puntuación por parte.}\label{tab:puntos-parte}
\begin{tabular}{@{}lcc@{}}
\toprule
Parte & Puntos & Nota obtenida \\
\midrule
A. Cuestiones (20 de 2 puntos) & 40 & \\
B. Problemas (6 de 10 puntos) & 60 & \\
\midrule
\textbf{Examen} & \textbf{100} & \\
\midrule
C. Proyecto de fase (se califica aparte) & 100 & \\
\bottomrule
\end{tabular}
\end{table}

\begin{table}[h]
\centering\small
\caption{Puntuación por módulo.}\label{tab:puntos-modulo}
\begin{tabularx}{\textwidth}{@{}l>{\raggedright\arraybackslash}X>{\raggedright\arraybackslash}Xcc@{}}
\toprule
Módulo & Cuestiones & Problemas & Puntos & Obtenidos \\
\midrule
M16. Preentrenamiento a escala & 2, 3, 18 (6~p) & B.1 y B.6 (20~p) & 26 & \\
M17. Hardware y sistemas & 1, 4, 5, 8, 14, 20 (12~p) & B.2 y B.3 (20~p) & 32 & \\
M18. Arquitecturas modernas & 9, 10, 11, 12, 13 (10~p) & B.4 (10~p) & 20 & \\
M19. Inferencia & 6, 7, 15, 16, 17, 19 (12~p) & B.5 (10~p) & 22 & \\
\midrule
\textbf{Total} & 20 cuestiones (40~p) & 6 problemas (60~p) & \textbf{100} & \\
\bottomrule
\end{tabularx}
\end{table}

El \cref{eje:fermi} combina los módulos 16 y 17; se asigna al~17 porque la mayor parte de su puntuación (apartados b, c y d) corresponde a la duración, la energía y la MFU.

\section{Umbral de aprobado}

\begin{itemize}[label=\textendash]
  \item \textbf{Examen de hito}: 70 puntos de 100 (70\,\%) o más.
  \item \textbf{Proyecto de fase}: 70 puntos de 100 o más en la rúbrica del \cref{tab:rubrica}.
  \item \textbf{Puerta de salida de la fase~IV}: además del examen, cada uno de los módulos M16 a M19 exige un test de módulo con al menos 9 respuestas correctas de 12 y una autoevaluación del laboratorio interactivo de la fase de al menos el 70\,\%.
\end{itemize}

Como orientación por módulo, una puntuación inferior al 70\,\% de los puntos de un módulo en el \cref{tab:puntos-modulo} (por ejemplo, menos de 18 de 26 en M16) indica que conviene repasarlo aunque el examen esté aprobado.

\section{Qué hacer si no se alcanza el umbral}

\begin{enumerate}[label=\arabic*.]
  \item Anotar cada fallo en el registro de errores clasificándolo como error de concepto, de cálculo o despiste.
  \item Repasar las secciones del \cref{tab:repaso} correspondientes a los fallos, empezando por los errores de concepto.
  \item Rehacer sin mirar la solución los ejercicios de nivel $\star\star$ y $\star\star\star$ de esos temas en los apuntes del módulo.
  \item Repetir el test del módulo afectado hasta alcanzar 9 de 12 y volver a presentarse al examen de hito con una versión nueva, no antes de una semana.
\end{enumerate}

\begin{table}[h]
\centering\small
\caption{Secciones de repaso según el fallo.}\label{tab:repaso}
\begin{tabularx}{\textwidth}{@{}>{\raggedright\arraybackslash}p{4.2cm}X@{}}
\toprule
Fallo & Repasar \\
\midrule
Cuestiones 1, 2, 3 & M17, §1.2 (regla $6ND$); M16, §4.5 (regla de 20 tokens) y §5.2 (Kaplan frente a Chinchilla) \\
Problema B.1 & M16, §4.2, §4.3 (forma cerrada), §4.5 y §4.6 (el mínimo es plano) \\
Problema B.2 & M17, §1.2, §2.5 (MFU), §6.1 y §6.2; M16, §5.3 (coste de la inferencia) \\
Cuestiones 4, 8; problema B.3 & M17, §5.1 (all-reduce en anillo), §5.3--§5.5 (paralelismo de datos, tensores y etapas; burbuja) y §5.6 (ZeRO) \\
Cuestiones 5, 20 & M17, §2.2 (intensidad crítica y roofline) y §2.5 (MFU) \\
Cuestión 14 & M17, §3.3 (softmax en línea) y §3.4 (FlashAttention) \\
Cuestiones 9, 10, 13; problema B.4 & M18, §4.1 (enrutado top-$k$), §4.2 (parámetros totales y activos), §4.3 (balanceo), §4.4 (capacidad) y §4.6 (casos públicos) \\
Cuestiones 7, 11 & M18, §2.1 y §2.3 (caché KV, MQA y GQA); M19, §2.2 \\
Cuestión 12 & M18, §3.1 y §3.2 (ventana deslizante y campo receptivo) \\
Cuestiones 6, 15; problema B.5 (a)--(d) & M19, §2.1 a §2.4 (\emph{prefill}, caché KV, intensidad y lote) \\
Cuestión 16; problema B.5 (e) & M19, §3.2 a §3.4 (decodificación especulativa) \\
Cuestión 17; problema B.5 (e) & M19, §4.3 (granularidad y coste de las escalas) \\
Cuestión 19 & M19, §6.2 y §6.3 (\emph{batching} continuo y PagedAttention) \\
Cuestión 18; problema B.6 & M16, §2.5 (LSH por bandas), §3.1--§3.2 (leyes de potencia) y §6.1--§6.2 (emergencia y métricas) \\
\bottomrule
\end{tabularx}
\end{table}

\paragraph{Siguiente paso.} Superado el hito, se continúa con la fase~V, módulo~20: ajuste fino.
